\chapter{Documentation API et Code}
\label{annexe:api}

\subsection*{Extrait du schéma SQL PostgreSQL}
\label{annexe:sql}

\begin{lstlisting}[language=SQL,caption={Extrait de la création des tables principales (identifiants UUID v4)}]
CREATE EXTENSION IF NOT EXISTS "uuid-ossp";

CREATE TABLE users (
  id                   UUID          PRIMARY KEY DEFAULT uuid_generate_v4(),
  phone                VARCHAR(20)   UNIQUE NOT NULL,
  email                VARCHAR(120)  UNIQUE,
  first_name           VARCHAR(80)   NOT NULL,
  last_name            VARCHAR(80)   NOT NULL,
  date_of_birth        DATE,
  password             VARCHAR(120)  NOT NULL,
  role                 VARCHAR(20)   NOT NULL DEFAULT 'passenger',
  verification_status  VARCHAR(20)   NOT NULL DEFAULT 'pending',
  is_verified          BOOLEAN       NOT NULL DEFAULT false,
  wilaya_code          INT,
  created_at           TIMESTAMPTZ   NOT NULL DEFAULT now()
);

CREATE TABLE trips (
  id                UUID           PRIMARY KEY DEFAULT uuid_generate_v4(),
  driver_id         UUID           NOT NULL REFERENCES users(id),
  vehicle_id        UUID           NOT NULL REFERENCES vehicles(id),
  from_wilaya_id    INT            NOT NULL REFERENCES wilayas(id),
  to_wilaya_id      INT            NOT NULL REFERENCES wilayas(id),
  from_city         VARCHAR(120)   NOT NULL,
  to_city           VARCHAR(120)   NOT NULL,
  departure_date    DATE           NOT NULL,
  departure_time    TIME           NOT NULL,
  price_per_seat    NUMERIC(10,2)  NOT NULL,
  available_seats   SMALLINT       NOT NULL,
  total_seats       SMALLINT       NOT NULL,
  status            VARCHAR(20)    NOT NULL DEFAULT 'active',
  created_at        TIMESTAMPTZ    NOT NULL DEFAULT now()
);

CREATE INDEX idx_trips_route_date ON trips
  (from_wilaya_id, to_wilaya_id, departure_date);
\end{lstlisting}



\subsection*{Extrait de code --- Factory de paiement}
\label{annexe:factory}

\begin{lstlisting}[language=Java,caption={Implémentation Node.js du patron \textit{Strategy/Factory} (extrait de \texttt{backend/src/services/patterns/payment/payment.factory.js})}]
const CIBStrategy      = require('./cib.strategy');
const EdahabiaStrategy = require('./edahabia.strategy');
const CashStrategy     = require('./cash.strategy');

class PaymentFactory {
  constructor(config = {}) {
    this.strategies = new Map();
    this.registerStrategy('cib',      new CIBStrategy(config.cib));
    this.registerStrategy('edahabia', new EdahabiaStrategy(config.edahabia));
    this.registerStrategy('cash',     new CashStrategy(config.cash));
  }

  registerStrategy(method, strategy) {
    this.strategies.set(method.toLowerCase(), strategy);
  }

  getStrategy(method) {
    const strategy = this.strategies.get(method.toLowerCase());
    if (!strategy) throw new Error(`Payment method not supported: ${method}`);
    return strategy;
  }

  async processPayment(method, paymentData) {
    return await this.getStrategy(method).process(paymentData);
  }

  async refundPayment(method, refundData) {
    return await this.getStrategy(method).refund(refundData);
  }
}

module.exports = PaymentFactory;
\end{lstlisting}

L'ajout d'une nouvelle méthode de paiement (p.~ex.~Baridi~Mob) consiste à créer une classe \texttt{BaridiStrategy} implémentant l'interface \texttt{process()} / \texttt{refund()}, puis à l'enregistrer dans la factory via \texttt{registerStrategy('baridi', ...)}~: aucun code existant n'est modifié, conformément au principe ouvert/fermé (OCP).

